% from main.tex
%%%%%%%%%%  scan idx 106  |  folio 93  %%%%%%%%%%
épimorphisme. Cela termine la démonstration de (1.4.1).

\textbf{Lemme 1.4.2.} \emph{Soient \(\mathcal C\) une catégorie abélienne, \(u:E\to F\) un morphisme de complexes de \(\mathcal C\) (resp.\ une flèche de \(\Dpx(\mathcal C)\)), \(G\) le cône de \(u\). Conditions équivalentes pour \(n\in\mathbb Z\) donné} :
\begin{enumerate}[label=(\roman*)]
\item \emph{\(\Ho^i(G)=0\) pour \(i\geq n\)}~;
\item \emph{\(\Ho^i(u)\) est un isomorphisme pour \(i>n\) et un épimorphisme pour \(i=n\)}.
\end{enumerate}
La preuve est immédiate à partir de la suite exacte de cohomologie.

\textbf{Définition 1.4.3.} \emph{Nous dirons que \(u\) est un \(n\)-quasi-isomorphisme (resp.\ un \(n\)-isomorphisme) si \(u\) vérifie les conditions équivalentes de (1.4.2).}

Cela étant, on déduit aussitôt de (1.4.1) le

\textbf{Corollaire 1.4.4.} \emph{Soient \(S,S,\mathcal C,\mathcal C_0\) comme en (1.4.1). Soit \(u:E\to F\) un \((n+1)\)-quasi-isomorphisme de complexes de \(\mathcal C_S\) (\(n\in\mathbb Z\) donné) tel que \(\Ho^n(G)\) soit de \(\mathcal C_0\)-type fini. Alors l'ensemble des \(X\) au-dessus de \(S\) tels qu'il existe \((M,r,s)\) comme en (1.4.1) tels que \(u'\) soit un \(n\)-quasi-isomorphisme est un raffinement de \(S\).}

De manière imagée, on peut, moyennant la condition de finitude sur \(\Ho^n(G)\), localement «~gagner un cran~», i.e. modifier localement \(u\) en degré \(n\) par l'addition à \(E^n\) d'un objet de \(\mathcal C_0\) de manière à faire de \(u\) un \(n\)-quasi-isomorphisme.

\textbf{Corollaire 1.4.5.} Soient \(\mathcal C\) une catégorie abélienne, \(n\in\mathbb Z\), et \(u:E\to F\) un \(n\)-quasi-isomorphisme de complexes de \(\mathcal C\). Il existe un diagramme commutatif de \(\Cpx(\mathcal C)\)
\[
\begin{tikzcd}
E \arrow[r,"v"] \arrow[d,"u"'] & E' \arrow[dl,"u'"]\\
F &
\end{tikzcd}
\]
